% Part: first-order-logic
% Chapter: model-theory
% Section: partial-iso

\documentclass[../../../include/open-logic-section]{subfiles}

\begin{document}

\olfileid{mod}{bas}{dlo}
\section{گھنیاں خطی ترتیباں}

\begin{defn}
  \emph{سِراں توں خالی گھنی خطی ترتیب}
  (dense linear ordering without endpoints) اوہ !!a{structure}
  $\Struct{M}$ اے جیہڑی اکّو دو-تھانواں والا
  !!{predicate}~$<$ رکھن والی !!{language} دی اے تے تھلے دتے
  جملے پورے کردی اے:
  \begin{enumerate}
  \item $\lforall[x][\lnot x < x]$؛
  \item $\lforall[x][\lforall[y][\lforall[z][(x < y \lif (y < z \lif x
    <z ))]]]$؛
  \item $\lforall[x][\lforall[y][(x< y \lor \eq[x][y] \lor y < x)]]$؛
  \item $\lforall[x][\lexists[y][x < y]]$؛
  \item $\lforall[x][\lexists[y][y < x]]$؛
  \item $\lforall[x][\lforall[y][(x < y \lif \lexists[z][(x < z \land
        z < y))]]]$۔
 \end{enumerate}
\end{defn}

\begin{thm}\ollabel{thm:cantorQ}
  سِراں توں خالی ہر دو !!{enumerable} گھنیاں خطی ترتیباں
  ہم شکل ہوندیاں نیں۔
\end{thm}

\begin{proof}
  $\Struct{M_1}$ تے $\Struct{M_2}$ سِراں توں خالی
  !!{enumerable} گھنیاں خطی ترتیباں ہون، جتھے
  ${<_1} = \Assign{<}{M_1}$ تے ${<_2} =
  \Assign{<}{M_2}$؛ انہاں دے وچکار ساریاں جزوی ہم شکلیاں دا
  سیٹ $\PIso{I}$ ہووے۔ $\PIso{I}$ غیر خالی اے کیوں جے گھٹ توں گھٹ
  $\emptyset \in \PIso{I}$۔ اسیں دکھاندے آں کہ $\PIso{I}$،
  اگے-پیچھے خاصیت پوری کردا اے۔ تاں
  $\Struct{M_1} \iso[p] \Struct{M_2}$، تے
  \olref[pis]{thm:p-isom1} توں قضیہ نکلدا اے۔

  $\PIso{I}$ دی اگے والی خاصیت دکھاؤن لئی $p \in
  \PIso{I}$ لو۔ جے نقشہ خالی ہووے، تاں دوجی ساخت دے غیر خالی دائرے
  وچوں کوئی وی عنصر چن کے دتے عنصر نال اوہدا جوڑا شامل کر لو۔ جے دتا
  عنصر پہلے ای نقشے دے دائرے وچ اے، تاں نقشہ آپ ہی مطلوبہ توسیع اے۔
  اگے غیر خالی نقشے تے دائرے توں باہر دتے عنصر دا حال اے؛ لسٹ وچ
  نقشے دے دائرے دے سارے عناصر لو۔ $p(a_i) = b_i$، $i = 1$، \dots،~$n$ لئی
  ہووے؛ دائرے دی لسٹ نوں ترتیب دے کے فرض کرو $a_1 <_1 a_2 <_1 \cdots <_1
  a_n$۔ کوئی $a \in \Domain{M_1}$ دتا ہووے؛ $b \in \Domain{M_2}$
  ایویں لبھو:
  \begin{enumerate}
  \item جے $a <_1 a_1$، تاں $b \in \Domain{M_2}$ ایسا لو
    کہ $b <_2
    b_1$؛
  \item جے $a_n <_1 a$، تاں $b \in \Domain{M_2}$ ایسا لو
    کہ $b_n <_2 b$؛
 \item جے $a_i <_1 a <_1 a_{i+1}$ کسے $i$ لئی، تاں
   $b \in
   \Domain{M_2}$ ایسا لو کہ $b_i <_2 b <_2 b_{i+1}$۔
  \end{enumerate}
  مطلوبہ خاصیت والا $b$ ہمیشہ لبھ جاندا اے، کیوں جے
  $\Struct{M_2}$ سِراں توں خالی گھنی خطی ترتیب اے۔
  $q = p \cup \{ \langle a, b \rangle \}$ مقرر کرو؛ ایویں
  $q \in \PIso{I}$، $p$ دا مطلوبہ ودھاؤ اے۔ ایس نال اگے
  والی خاصیت ثابت ہوندی اے۔ پیچھے والی خاصیت دا ثبوت ایہو
  جیہا اے۔ ایس لئی $\Struct{M_1} \iso[p]
  \Struct{M_2}$؛ \olref[pis]{thm:p-isom1} نال
  $\Struct{M_1} \iso
  \Struct{M_2}$۔
\end{proof}

\begin{prob}
  \olref[mod][bas][dlo]{thm:cantorQ} دا ثبوت ایہ جانچ کے پورا
  کرو کہ $\PIso{I}$ پیچھے والی خاصیت پوری کردا اے۔
\end{prob}

\begin{rem}
  $\Struct{S}$ کوئی وی سِراں توں خالی !!{enumerable} گھنی
  خطی ترتیب ہووے۔ تاں (\olref{thm:cantorQ} مطابق)
  $\Struct{S} \iso
  \Struct{Q}$، جتھے $\Struct{Q} = (\Rat, <)$ اوہ
  !!{enumerable} گھنی خطی ترتیب اے جیہدا دائرہ ناطق اعداد دا
  سیٹ $\Rat$ اے۔ ہن \olref[thm]{remark:R} وچ آئی
  !!{structure}~$\Struct{R} =
  (\Real, <)$ نوں فیر ویکھو۔ اسیں ویکھیا سی کہ اک
  !!a{enumerable} !!{structure}~$\Struct{S}$ اے جیہدے لئی
  $\Struct{R}
  \elemequiv \Struct{S}$۔ پر $\Struct{S}$ سِراں توں خالی
  !!a{enumerable} گھنی خطی ترتیب اے، ایس لئی اوہ
  !!{structure}~$\Struct{Q}$ دے نال ہم شکل (تے ایس سبب
  جملہ وار ہم ارز) اے۔ جملہ وار ہم ارزی دی تعدی خاصیت توں
  $\Struct{R} \elemequiv
  \Struct{Q}$۔ (ایہ نتیجہ اسیں اوہو اگے-پیچھے استدلال نال
  $\Struct{R} \iso[p] \Struct{Q}$ قائم کرکے براہِ راست وی
  دکھا سکدے ساں۔)
\end{rem}
\end{document}
